% =====================================================================
\section{Experimental setup}\label{sec:setup}
% =====================================================================
\subsection{Implementation details}
All experiments were implemented in PyTorch \citep{paszke2019pytorch} with the \texttt{timm} model library
\citep{wightman2019timm} and scikit-learn \citep{pedregosa2011scikit}, and run on a single \res{gpu} GPU. Mixed
precision (fp16) \citep{micikevicius2018mixed} was used for training and inference, and Grad-CAM maps were computed
with the \texttt{pytorch-grad-cam} library \citep{gildenblat2021gradcam}. Random seeds were fixed to 42 for
data splitting and pretraining and to $42+k$ for fold $k$. Table~\ref{tab:hparams} lists all hyper-parameters; none of
them was tuned on the test set. After early stopping, the number of epochs actually run in the five folds was
(\res{epochsrun}).

\begin{table}[!htbp]
\centering\small
\caption{Training and inference hyper-parameters. Values marked $\dagger$ are shared by both stages.}
\label{tab:hparams}
\begin{adjustbox}{max width=\linewidth}
\begin{tabular}{@{}lll@{}}
\toprule
Setting & Stage 1 (APTOS-2019) & Stage 2 (IDRiD, per fold)\\
\midrule
Initialisation & ImageNet-21k backbone & stage-1 weights\\
Epochs (max.) & 10 & 25, early stopping (patience 8)\\
Peak learning rate (head / backbone) & $3\times10^{-4}$ / $1\times10^{-4}$ & $2\times10^{-4}$ / $6.7\times10^{-5}$\\
EMA decay & 0.999 & 0.99\\
Optimiser$^\dagger$ & \multicolumn{2}{l}{AdamW \citep{loshchilov2019decoupled}, weight decay $10^{-4}$}\\
Schedule$^\dagger$ & \multicolumn{2}{l}{one-cycle cosine, 10\% warm-up \citep{smith2019super}}\\
Batch size$^\dagger$ & \multicolumn{2}{l}{16; gradient clipping at $\ell_2$ norm 2.0}\\
Input size$^\dagger$ & \multicolumn{2}{l}{$448\times448$, Ben Graham pre-processing}\\
Loss weights$^\dagger$ & \multicolumn{2}{l}{$\lambda_r=0.5$ (ordinal), $\lambda_m=0.01$ (sparsity), label smoothing 0.05}\\
Class weights$^\dagger$ & \multicolumn{2}{l}{$w_c\propto n_c^{-1/2}$, normalised to mean 1, computed on training data}\\
Checkpoint selection$^\dagger$ & \multicolumn{2}{l}{max.\ validation accuracy + QWK (EMA weights, fixed cut-points, no TTA)}\\
Inference & -- & 5 fold models $\times$ 4 flips = 20 predictions per image\\
\bottomrule
\end{tabular}
\end{adjustbox}
\end{table}

\subsection{Evaluation metrics}
Let $y_i$ and $\hat y_i$ be the reference and predicted grades of test image $i$, $i=1,\dots,n$, and let $O$ be the
$5\times5$ confusion matrix with entries $O_{jk}$ counting images of grade $j$ predicted as grade $k$.

\paragraph{Five-class grading.} \emph{Accuracy} is the fraction of exactly correct grades. Because the classes are
imbalanced, we also report \emph{balanced accuracy}, the mean of the per-class recalls, and the macro-averaged
precision, recall and F1, which weight every grade equally. The \emph{quadratic weighted kappa}
\citep{cohen1968weighted} measures agreement beyond chance while penalising errors by the squared grade distance:
\begin{equation}
\kappa_w = 1 - \frac{\sum_{j,k} W_{jk}\,O_{jk}}{\sum_{j,k} W_{jk}\,E_{jk}},\qquad
W_{jk} = \frac{(j-k)^2}{(K-1)^2},
\end{equation}
where $E$ is the expected confusion matrix under independence (outer product of the marginal histograms, scaled to
$n$) and $K=5$. We also report unweighted Cohen's kappa, the Matthews correlation coefficient (MCC), the
\emph{within-one-grade accuracy} $\frac1n\sum_i \mathbb{1}[|y_i-\hat y_i|\le1]$, which quantifies how many errors are
between adjacent grades, and macro one-vs-rest ROC AUC and average precision computed from the ensemble
probabilities.

\paragraph{Referable DR.} With TP, FN, FP and TN counted for the binary task (grade $\geq2$ vs.\ $<2$), we report
sensitivity $\mathrm{TP}/(\mathrm{TP}+\mathrm{FN})$, specificity $\mathrm{TN}/(\mathrm{TN}+\mathrm{FP})$, positive and
negative predictive values, accuracy, F1 and the threshold-free ROC AUC of $P(\text{referable})$. Sensitivity is the
clinically most important quantity, because a missed referable case delays treatment.

\subsection{Statistical analysis}
\paragraph{Confidence intervals.} For each test metric we compute a 95\% percentile bootstrap interval
\citep{efron1993introduction}: the $n$ test images are resampled with replacement 2000 times, the metric is recomputed
on each resample, and the 2.5th and 97.5th percentiles are reported. Predictions are fixed during resampling, so the
interval reflects uncertainty due to the finite test set, not training randomness.

\paragraph{Paired comparisons.} All configurations are evaluated on the same test images, so we compare each
ablation with the full model using the two-sided exact McNemar test on per-image correctness
\citep{dietterich1998approximate}. If $b$ images are correct only for the full model and $c$ only for the ablation,
the $p$-value is
\begin{equation}
p = \min\!\Big(1,\; 2\sum_{i=0}^{\min(b,c)} \binom{b+c}{i}\,2^{-(b+c)}\Big).
\end{equation}
We report $p$-values without claiming significance for a single threshold; seven ablations are compared, so a
Bonferroni-adjusted level of $0.05/7\approx0.007$ is the conservative reference.

\paragraph{Cross-validation estimates.} We additionally report OOF metrics and their mean and standard deviation across
folds. These are slightly optimistic, because each fold's checkpoint is selected on the same fold
\citep{varma2006bias}; they are given for completeness and are not used as the main result.

\subsection{Compared configurations}
All configurations use identical splits, seeds, augmentation, schedule and inference, and differ in one factor only:
\begin{itemize}
\item \textbf{Plain}: the same backbone and $1\times1$ projection, followed by global average pooling and the same
  dual head and loss. This is the standard transfer-learning baseline and isolates the effect of the whole
  cross-expert head.
\item \textbf{w/o APTOS}: no stage-1 pretraining; the model is initialised from ImageNet-21k only.
\item \textbf{CLAHE}: CLAHE + bilateral filtering instead of Ben Graham pre-processing, with stage 1 repeated.
\item \textbf{w/o ordinal head}: $\lambda_r=0$; the regression output is untrained, and the decision-rule selection on OOF
  data (Section~\ref{sec:decision}) decides whether the score $s$ or the argmax of $p$ is used.
\item \textbf{w/o lesion map}: $m\equiv1$; the lesion expert receives all tokens with equal weight.
\item \textbf{w/o cross-attention}: the experts are only normalised, $G'=\mathrm{LN}(G)$, $L'=\mathrm{LN}(L)$.
\item \textbf{w/o gate}: the gate is fixed to $g=0.5$, i.e.\ equal contribution of both experts.
\end{itemize}
For the architectural variants, stage 1 is repeated so that pretraining and fine-tuning use the same architecture.

\subsection{Reproducibility}
The notebook, the export cell that writes every reported number and figure, the ablation variants and the script that
inserts the numbers into this manuscript are released together (see Code availability). The held-out test indices and
per-image predictions are stored in the exported results, so every table can be regenerated and every comparison
re-tested without retraining.
\sectionend{setup}
